The main observation is that
\[
2197=13^3,\qquad \phi(2197)=13^3-13^2=2028.
\]
So a first return to \(1\) at exponent \(2028\) asks for \(m\) to be a primitive root modulo \(2197\).

Primitive roots exist for odd prime powers.  Hence the number in one reduced residue system is
\[
\phi(\phi(2197))=\phi(2028).
\]
Since
\[
2028=2^2\cdot3\cdot13^2,
\]
this is
\[
2028\cdot\frac12\cdot\frac23\cdot\frac{12}{13}=624.
\]

The starting slot is technically chosen from \(0,\ldots,9999\), not from a single complete period modulo \(2197\).  This creates a boundary issue because
\[
10000=4\cdot2197+1212.
\]
Thus one should count primitive roots in the first \(1212\) residue classes.  I will ignore that boundary effect and use the single-period probability, since the modulus is the intended main structure:
\[
P=\frac{624}{2197}.
\]
Reducing by \(13\),
\[
P=\frac{48}{169}.
\]
Therefore
\[
p+q=48+169=217.
\]
So my answer is
\[
\boxed{217}.
\]
